% !TeX program = lualatex
% !TeX encoding = UTF-8 Unicode
%---------------------------------------------------------------------------
% One-page academic CV, English. Template from the docs-resume skill; every
% name, place, date, title and URL below is a placeholder.
%
% One page by construction, not by shrinking a two-page layout: extarticle at
% 9pt, 1.3 cm side margins, Libertinus Serif (it holds up at this size), each
% publication one hanging paragraph (authors, title, venue) instead of a title
% line plus an author line, education detail at body size, and honors beside
% service in two columns at 7.6pt. Colour only on links. Section heads are
% letterspaced small caps over a 0.4pt hair rule.
%
% Cross-references: every publication carries \hypertarget{pub:<n>}. A bullet
% written \result[<n>]{...} and a method name written \pubref{<n>}{Name} jump
% to it. Numbers run newest highest, so adding a paper never shifts an id.
%
% Build: lualatex (fontspec). Overleaf: Menu > Compiler = LuaLaTeX.
% The Chinese sibling is cv-onepage-zh_CN.tex, the bundle cv-onepage-bilingual.tex.
%---------------------------------------------------------------------------
\documentclass[a4paper,9pt]{extarticle}

\usepackage{fontspec}
\setmainfont{LibertinusSerif-Regular.otf}[
  BoldFont=LibertinusSerif-Bold.otf,ItalicFont=LibertinusSerif-Italic.otf,
  BoldItalicFont=LibertinusSerif-BoldItalic.otf,Numbers=Proportional]
\newfontfamily\headface{LibertinusSerif-Regular.otf}[
  BoldFont=LibertinusSerif-Bold.otf,Letters=SmallCaps,LetterSpace=8]
\usepackage{unicode-math}
% By file name: the family name is not in every system font database.
\setmathfont{LibertinusMath-Regular.otf}
\usepackage[verbose=silent]{microtype}
\usepackage[left=1.3cm,right=1.3cm,top=1.0cm,bottom=1.0cm]{geometry}
\usepackage{xcolor}
\usepackage{titlesec}
\usepackage{enumitem}
\usepackage{xurl}
\usepackage[hidelinks]{hyperref}

\definecolor{ink}{HTML}{1A1A1A}
\definecolor{grey}{HTML}{6B6B6B}
\definecolor{hair}{HTML}{BDBDBD}
\definecolor{link}{HTML}{1F3A5F}   % dark navy: reads near-black, blue on inspection
\hypersetup{colorlinks=true,urlcolor=link,linkcolor=link,pdfborder={0 0 0},
  pdftitle={Firstname Lastname, Curriculum Vitae},pdfauthor={Firstname Lastname}}
\urlstyle{same}
\pagestyle{empty}
\color{ink}
\raggedright
\linespread{1.15}
\setlength{\parindent}{0pt}
\setlength{\tabcolsep}{0pt}
\setlength{\emergencystretch}{1.5em}

% Section head: letterspaced small caps, a hair rule about 3pt under the
% baseline. A larger negative \vspace pulls the rule into the descenders.
\titleformat{\section}
  {\headface\bfseries\normalsize}{}{0pt}{}
  [\vspace{-1.0pt}{\color{hair}\titlerule[0.4pt]}]
\titlespacing*{\section}{0pt}{5.5pt plus 1pt minus 1pt}{2.5pt}

% One line, left and right halves: dates and places align right, no dashes.
\newcommand{\headline}[2]{\par\noindent\hbox to \linewidth{#1\hfil #2}\par}

% Institution or employer entry, two lines.
% #1 name (linked)  #2 place  #3 dates  #4 detail line (degree, advisors / role, mentors)
\newcommand{\entry}[4]{\addvspace{2.5pt}%
  \headline{{\bfseries #1}}{{\footnotesize\color{grey}#2\hspace{1.1em}#3}}%
  \nopagebreak{\normalsize\upshape\color{ink}#4}\par\addvspace{1pt}}
\newcommand{\sidesize}{\fontsize{7.6}{9.2}\selectfont}

% Research bullets: a bold heading, a grey tag, a jump to the paper, one paragraph.
\setlist[itemize]{leftmargin=1.0em,label={\color{grey}\textbullet},
  topsep=0.8pt,itemsep=2.5pt,parsep=0pt,partopsep=0pt}
% #1 publication number (optional)  #2 heading  #3 tag  #4 body
\newcommand{\result}[4][]{\item
  {\bfseries #2}%
  \if\relax\detokenize{#3}\relax\else\ {\footnotesize\color{grey}[#3]}\fi%
  \if\relax\detokenize{#1}\relax\else\ {\footnotesize\color{grey}\hyperlink{pub:#1}{[#1]}}\fi
  \ #4}
% A method name in the opening that jumps to its entry: [\pubref{10}{MethodA}]
\newcommand{\pubref}[2]{\hyperlink{pub:#1}{#2}}

% Publications: one hanging paragraph each. Venue in an \mbox; \nolinebreak keeps
% [link] and the badge from wrapping onto a line of their own.
% #1 badge after the link (optional)  #2 number  #3 title  #4 venue  #5 authors  #6 url
\newlist{publist}{itemize}{1}
\setlist[publist]{leftmargin=2.0em,labelwidth=1.65em,labelsep=0.35em,align=left,label={},
  topsep=1.5pt,itemsep=1.8pt,parsep=0pt,partopsep=0pt}
\newcommand{\pub}[6][]{%
  \item[{\color{grey}\hypertarget{pub:#2}{[#2]}}]#5.
  \if\relax\detokenize{#6}\relax #3\else\href{#6}{\textcolor{ink}{#3}}\fi.
  \mbox{\bfseries\color{grey}#4}.%
  \if\relax\detokenize{#6}\relax\else\nolinebreak\ \href{#6}{[link]}\fi
  \if\relax\detokenize{#1}\relax\else\nolinebreak\ \mbox{#1}\fi}

% One-line rows with the year at the right, and label/value rows. The label
% column is measured from the WIDEST label (see \settowidth below).
\newcommand{\row}[2]{\headline{#1}{{\footnotesize\color{grey}#2}}\addvspace{1.5pt}}
\newlength{\kvlab}
\newcommand{\kv}[2]{\noindent
  \makebox[\kvlab][l]{\bfseries #1}%
  \begin{minipage}[t]{\dimexpr\linewidth-\kvlab\relax}\raggedright\strut #2\end{minipage}\par\addvspace{1.5pt}}

\begin{document}
\settowidth{\kvlab}{\bfseries Languages}\addtolength{\kvlab}{0.9em}

%---- Name, status, one contact line, one hairline
\headline{{\fontsize{19}{21}\selectfont\bfseries Firstname Lastname}}%
  {{\small\color{grey}Ph.D. candidate, University A\quad|\quad Available from Month YYYY}}
\vspace{1pt}
{\small\color{grey}
\href{mailto:firstname.lastname@example.org}{firstname.lastname@example.org}\quad|\quad
\href{https://example.org/scholar}{Google Scholar}\quad|\quad
(+1) 555-0100\quad|\quad City, Country\par}
\vspace{3pt}
{\color{hair}\hrule height 0.4pt}
\vspace{4pt}

%---- Research focus: a bold run-in label, one paragraph, five lines
{\bfseries Research focus}\quad
My research makes Model Class Y reliable under Condition X, along two lines of work. On the data side, I develop methods that decide what a model learns from, for Capability One [\pubref{5}{MethodA}] and Capability Two [\pubref{9}{MethodC}], so that every labelled example goes further and unlabelled audio becomes usable for training. On the model side, I build training objectives [\pubref{8}{MethodB}] that keep predictions stable when the recording conditions change, so one model serves every device. My broader vision is Vision V, stated in a phrase the field already uses, which the two lines approach from opposite ends.

%---- Education: two lines per degree
\section{Education}
\entry{\href{https://example.org/university-a}{\textcolor{ink}{University A}}}{City, Country}{YYYY.MM -- YYYY.MM}
  {Ph.D. in Computer Science, expected Month YYYY, advised by \href{https://example.org/advisor-one}{Prof.\ Advisor One}. Internships at \href{https://example.org/company-b}{Company B} and \href{https://example.org/company-c}{Company C}}
\entry{\href{https://example.org/university-b}{\textcolor{ink}{University B}}}{City, Country}{YYYY.MM -- YYYY.MM}
  {B.Sc. in Electrical Engineering, research advised by \href{https://example.org/advisor-two}{Prof.\ Advisor Two}}

%---- Experience: only what the target reader weighs most
\section{Research Experience}
\entry{\href{https://example.org/company-b}{\textcolor{ink}{Team Name, Lab Name, Company B (now New Lab Name)}}}{City, Country}{YYYY.MM -- present}
  {Research Intern, mentored by \href{https://example.org/mentor-one}{Dr.\ Mentor One} and \href{https://example.org/mentor-two}{Dr.\ Mentor Two}}
\begin{itemize}
  \result[9]{Label-efficient speech recognition}{MethodC, Venue YYYY submission}{Proposed MethodC, in which each pseudo-label is weighted by how often two decoders agree on it, addressing the label noise that limits self-training in low-resource languages. An agreement score filters every utterance, and a schedule raises the threshold as the model improves. Across eight languages at two model scales, MethodC cut word error rate by about 12\% relative to the strongest baseline with 40\% fewer labelled hours.}
  \result{Audio-visual extension}{BenchmarkC}{Extended MethodC to audio-visual input, studying whether lip movements help most where the audio is noisiest and whether the agreement score carries over to video.}
\end{itemize}

\entry{\href{https://example.org/company-c}{\textcolor{ink}{Research Group, Company C}}}{City, Country}{YYYY.MM -- YYYY.MM}
  {Research Intern, mentored by \href{https://example.org/mentor-three}{Dr.\ Mentor Three}, \href{https://example.org/mentor-four}{Dr.\ Mentor Four}, and \href{https://example.org/mentor-five}{Prof.\ Mentor Five}}
\begin{itemize}
  \result[8]{Robust audio classification}{MethodB, ICML 2025 Spotlight}{Introduced MethodB, a training objective that matches a classifier's predictions across simulated recording conditions, so a model trained on studio audio holds up on phone and far-field recordings. The conditions are sampled from a learned room and device model rather than a fixed augmentation list. MethodB led all baselines on three shifted test sets, +5.3/+3.1 on BenchmarkA/BenchmarkB, at no extra inference cost.}
  \result[7]{Data valuation for audio corpora}{arXiv 2025}{Co-developed an estimator that scores each training clip by its effect on validation loss and removes the clips that hurt. The effect is approximated from the gradients of a few saved checkpoints, so scoring a million clips takes one GPU-day. On one public corpus it ranked most mislabelled clips in its lowest tenth, and dropping that tenth raised accuracy by 2.0 points over training on the full corpus.}
\end{itemize}

\entry{\href{https://example.org/company-d}{\textcolor{ink}{Applied Research Team, Company D}}}{City, Country}{YYYY.MM -- YYYY.MM}
  {Research Intern, mentored by \href{https://example.org/mentor-six}{Dr.\ Mentor Six}}
\begin{itemize}
  \result[4]{On-device keyword spotting}{arXiv 2023}{Proposed a keyword spotter that shares one small encoder across all keywords and adds a new keyword through a short enrolment step, in place of one model per keyword. The encoder is trained once with a metric-learning objective, and enrolment needs three spoken examples. It ran in under 1 MB of memory and stayed within one point of a model ten times its size.}
  \result[2]{Speaker verification under noise}{Interspeech 2022}{Introduced a noise-aware speaker embedding that down-weights the frames a small detector marks as corrupted, in place of denoising the audio first. The detector is trained on synthetic corruptions and needs no clean reference. It lowered the equal error rate by 15\% relative on the noisy trials of two benchmarks, with no loss on clean trials.}
\end{itemize}

%---- Publications: the complete list, newest highest
\section{Publications}
{\small\color{grey}9 publications, 5 first-author, 3 of them accepted (ICML, Interspeech $\times$2). Built the speech track of the Benchmark Title.}

{\small
\begin{publist}
\pub{9}{MethodC: Paper Title Nine, a Placeholder under Review}{arXiv 2026}{\textbf{F. Lastname}, C. Coauthor, D. Coauthor, E. Coauthor, G. Coauthor, A. Advisor}{https://example.org/scholar}
\pub{8}{MethodB: Paper Title Eight, a Placeholder for a Robust Classification Paper}{ICML 2025 Spotlight}{\textbf{F. Lastname}, H. Coauthor, A. Advisor}{https://example.org/paper-8}
\pub[\href{https://example.org/ranking}{[Top-5 of the Week]}]{7}{Paper Title Seven, a Placeholder for a Collaboration}{arXiv 2025}{G. Coauthor, D. Coauthor, \textbf{F. Lastname}, A. Advisor}{https://example.org/paper-7}
\pub{6}{Benchmark Title: A Placeholder Benchmark with Dozens of Authors}{TACL 2024}{{\bfseries\color{ink}Built the speech track.} K. Author et al.\ (24 authors)}{https://example.org/paper-6}
\pub{5}{MethodA: Paper Title Five, a Placeholder for Data Selection}{Interspeech 2023}{\textbf{F. Lastname}, I. Coauthor, A. Advisor}{https://example.org/paper-5}
\pub{4}{Paper Title Four, a Placeholder for a Short Study}{arXiv 2023}{\textbf{F. Lastname}, J. Coauthor, B. Advisor}{https://example.org/paper-4}
\pub{3}{Paper Title Three, a Placeholder for a Collaboration in a Second Field}{CVPR 2023}{L. Coauthor, M. Coauthor, N. Coauthor, \textbf{F. Lastname}, B. Advisor}{https://example.org/paper-3}
\pub{2}{Paper Title Two, a Placeholder for an Early First-Author Paper}{Interspeech 2022}{\textbf{F. Lastname}, O. Coauthor, B. Advisor}{https://example.org/paper-2}
\pub{1}{Paper Title One, a Placeholder for the Earliest Preprint}{arXiv 2021}{P. Coauthor, \textbf{F. Lastname}}{https://example.org/paper-1}
\end{publist}
}

%---- Honors beside service and skills, in smaller type
\vspace{2pt}
\noindent\begin{minipage}[t]{0.56\linewidth}
\section{Honors and Awards}
\sidesize
\row{\textbf{Fellowship A} (CUR 20,000), University A}{YYYY}
\row{\textbf{Best Paper Award}, Workshop Name at Venue}{YYYY}
\row{\textbf{Scholarship B} (top 5\%), University B}{YYYY--YYYY}
\row{\textbf{Travel Grant}, Conference Name}{YYYY}
\row{\textbf{Dean's List}, University B}{YYYY}
\end{minipage}\hfill
\begin{minipage}[t]{0.40\linewidth}
\section{Service and Skills}
\sidesize
\kv{Teaching}{TA, \href{https://example.org/course}{COURSE 101: Course Title}, YYYY and YYYY}
\kv{Service}{Reviewer for ICML (YYYY, YYYY), ICASSP (YYYY)}
\kv{Skills}{Python, C++, PyTorch, JAX, torchaudio, Slurm}
\kv{Languages}{Chinese (native), English (professional)}
\end{minipage}

\end{document}
